\subsection{引戻し成分の幾何的同定と局所模型の帰結}

三角形が共点の場合の基本1形式は
\[
 \eta=c_0\,dw-w\,dc_0+(f-wL)(x\,dy-y\,dx).
\]
これは $(x,y,w)$ だけに依存し、重み付き射影平面
$\mathbb P(1,1,3)$ 上の重み六の基本1形式である。
逆に、そのような1形式は一意に
\[
 c_3\,dw-w\,dc_3+(f_4-wL_1)(x\,dy-y\,dx)
\]
と書ける。パラメータ数は $4+5+2=11$ である。

一般の重み付き射影は、三次元の一次形式空間の二平面の選択と、
$w$ に加える三次式を二変数三次式で割った選択からなる。
そのパラメータ数は $2+(10-4)=8$ である。
従って引戻し族の次元は $8+11=19$。
一般点では、二平面は $c_3$ の三つの一次因子の張る空間から復元でき、
残る六つの射影パラメータは前補題の $z$ を含む三次移動から復元できるので、
この数え上げに正次元の余分なファイバーはない。

上の全局所模型は二つの被約枝を持つ。
$u=0,b=0$ は一般対数型の閉包であり、
$a=b=0$ がこの $\mathbb P(1,1,3)$ への射影による引戻しの枝である。
両枝の被約交叉は滑らかな十八次元である。
引戻し枝の一般点では $u$ が単元なので、その関数体を $K$ とすると横断環は
\[
 K[a,b]/(a^5,ab,b^2)
\]
となる。基底は $1,a,a^2,a^3,a^4,b$ であり、長さは六である。
従って引戻し成分は一般点で非被約である。

関数体 $K$ 上の表示には係数体の非標準的な選択を含める。
切片を得るために有限分離拡大を用いた場合にも、この環の同型型は $K$ 上へ下降する。
実際、極大イデアルのHilbert関数は $(1,2,1,1,1)$、socle次元は二であり、これらは下降する。
係数体を選び、$a^4\ne0$ となる $a$ と、socleの中で $a^4$ に独立な $b$ を取れば、
$ab=b^2=0$ かつ $1,a,a^2,a^3,a^4,b$ が基底となる。
ここで $a$ は無限体上の非空開条件から選べ、$b$ は $a^4$ 以外のsocle方向である。

\begin{remark}[二つのprimary成分と交叉の厚み]
\[
 (ab,b^2,a^5u)=(a^5,ab,b^2)\cap(b,u)
\]
である。従ってembedded primeはなく、第一のprimary成分が引戻し枝の厚みを表す。
このprimary成分と対数型枝のスキーム論的交叉は
$(b,u,a^5)$ であり、被約交叉の上に長さ五の厚みを持つ。
さらにHilbert--Burch行列
\[
 \begin{pmatrix}b&a^4u\\-a&0\\0&-b\end{pmatrix}
\]
の最大小行列式が三つの生成元を与えるので、滑らかな因子を除いた局所環は
一次元Cohen--Macaulayである。
一般の横断切断 $u=-a-\beta b$ の長さは七となる。
これらの交叉の厚みは、二つの被約成分の交叉次元だけからは読み取れない。
\end{remark}
